\textbf{Scale and scope.} One 2D environment with 2 modes, position-only observations, MLP heads, 10 seeds (5 in ablations), 200 rollouts per run. Many results are at the ceiling (success exactly $1.0$), so the study cannot discriminate between the GMM and DDPM heads on success; harder tasks (more modes, high-dimensional actions, images, contact) are needed for that and are not covered. \textbf{Baselines.} All three heads are our own reimplementations, with no per-method tuning; a tuned DDPM (e.g.\ a larger network, different schedule or step count, EMA) could change the DDPM rows, and we did not include IBC, BeT, IMLE or other published heads~\cite{florence2021implicit,shafiullah2022behavior,rana2025imle}. The DDPM MLP has one more hidden layer than the other heads. \textbf{Design choice.} The main start jitter (0.01) was selected from a pilot seed to be in the regime where MSE is partly working, and the main H1 verdict is conditional on it; the sweep shows MSE success from $0$ to near the mixtures depending on the jitter. \textbf{Mechanisms} (symmetry breaking via input-dependent mean, the timeouts of the chunked DDPM, the benefit of $M=1$) are hypotheses, not isolated by experiments. \textbf{Statistics.} Tests on 5--10 seeds with saturated, non-Gaussian metrics (several groups have zero variance) are indicative only; no correction for multiple comparisons. \textbf{Metrics.} Mode share is defined by the side at which the trajectory crosses $x=0.5$; the 40-step horizon makes timeouts count as failures.
